% Ad Familiares 10.21 (ad-familiares-10-21)
% Source: https://raw.githubusercontent.com/PerseusDL/canonical-latinLit/master/data/phi0474/phi056/phi0474.phi056.perseus-lat1.xml
% Fetched by scripts/fetch_latin.py

% Dateline (Perseus): Scr. in castris pnd. Id. Mai. a. 711 (43) .

\ciceroLetterOpener{PLANCVS CICERONI}

\ciceroSection{1} puderet me inconstantiae mearum litterarum, si non haec ex aliena levitate penderent. omnia feci qua re Lepido coniuncto ad rem p. defendendam minore sollicitudine vestra perditis resisterem ; omnia ei et petenti recepi et ultro pollicitus sum scripsique tibi biduo ante confidere me bono Lepido esse usurum communique consilio bellum administraturum. credidi chirographis eius, adfirmationi praesentis Laterensis, qui tum apud me erat reconciliaremque me Lepido fidemque haberem orabat. non licuit diutius bene de eo sperare ; illud certe cavi et cavebo ne mea credulitate rei p. summa fallatur.

\ciceroSection{2} Cum Isaram flumen uno die ponte effecto exercitum traduxissem pro magnitudine rei celeritatem adhibens, quod petierat per litteras ipse ut maturarem venire, praesto mihi fuit stator eius cum litteris, quibus ne venirem denuntiabat; se posse per se conficere negotium ; interea ad Isaram exspectarem. indicabo temerarium meum consilium tibi : nihilo minus ire decreram existimans eum socium gloriae vitare. putabam posse me nec de laude ieiuni hominis delibare quicquam et subesse tamen propinquis locis, ut, si durius aliquid esset, succurrere celeriter possem.

\ciceroSection{3} ego a non malus homo hoc suspicabar. at Laterensis, vir sanctis simus suo chirographo mittit mihi litteras nimisque desperans de se, de exercitu, de Lepidi fide querensque se destitutum, in quibus aperte denuntiat videam ne fallar ; suam fidem solutam esse ; rei p. ne desim. exemplar eius chirographi Titio misi ; ipsa chirographa omnia, et ea quibus credidi, et ea quibus fidem non habendam putavi, Laevo Cispio dabo perferenda, qui omnibus iis interfuit rebus.

\ciceroSection{4} accessit eo ut milites eius, cum Lepidus contionaretur, improbi per se, corrupti etiam per eos qui praesunt, Canidios Rufrenosque et ceteros quos cum opus erit scietis, conclamarent viri boni pacem se velle neque esse cum ullis pugnaturos duobus iam consulibus singularibus occisis, tot civibus pro patria amissis, hostibus denique omnibus iudicatis bonisque publicatis ; neque hoc aut vindicarat Lepidus aut sanarat.

\ciceroSection{5} hoc me venire et duobus exercitibus coniunctis obicere exercitum fidelissimum, auxilia maxima, principes Galliae, provinciam cunctam summae dementiae et temeritatis esse vidi, mihique, si ita oppressus essem remque publicam mecum prodidissem, mortuo non modo honorem sed misericordiam quoque defuturum. itaque rediturus sum nec tanta munera perditis hominibus dari posse sinam.

\ciceroSection{6} ut exercitum locis habeam opportunis, provinciam tuear, etiam si ille exercitus descierit, omniaque integra servem dabo operam, quoad exercitus hoc summittatis parique felicitate rem p. hic vindicetis. nec depugnare, si occasio tulerit, nec obsideri, si necesse fuerit, nec mori, si casus inciderit, pro vobis paratior fuit quisquam. qua re hortor te, mi Cicero, exercitum hoc traiciendum quam primum cures et matures, prius quam hostes magis conroborentur et nostri perturbentur. in quo si celeritas erit adhibita, res p. in possessione victoriae deletis sceleratis permanebit. fac valeas meque diligas.
